\section{后训练与 RL Infra：Agent 时代的核心系统}

访谈反复提到，Agent 范式很吃后训练。原因是 Agent 不再只回答单个问题，而是在环境中做多步行动。它需要 rollout 生成轨迹、reward/verifier 提供反馈、policy update 改变策略、evaluation 衡量端到端结果，还需要 scaffold 管理工具与任务。

\begin{figure}[H]
\centering
\includegraphics[width=0.94\textwidth]{agent-rl-infra-loop.png}
\caption{Agent post-train 与 RL infra 的闭环。}
\end{figure}

\begin{knowledgebox}{读图：Agent 后训练不是一个单点模块}
图中任务环境、rollout engine、reward/verifier、policy update、eval/cost 和 scaffold 构成闭环。它说明后训练不是“再微调一下”，而是一个工程系统：任务怎么产生、轨迹怎么采、反馈怎么定义、策略怎么更新、成本怎么衡量，都要同时解决。
\end{knowledgebox}

\subsection{Rollout 推理引擎到 Agent 核心系统}

描述里有一句关键话：系统从“以 rollout 推理引擎为核心”，转变为“以 Agent 为核心”的复杂系统。前者更像围绕模型生成答案与推理路径；后者需要整个环境闭环。对团队的要求也变化了：必须足够敏捷，能快速开发适配新范式的 RL infra。

\noindent\begin{tabularx}{\textwidth}{p{0.22\textwidth}p{0.34\textwidth}X}
\toprule
阶段 & 训练/评估对象 & 团队能力要求 \\
\midrule
Chat 后训练 & 回答质量、偏好、推理链、对齐 & 数据清洗、偏好标注、指令/偏好优化。 \\
Reasoning 后训练 & 可验证数学、代码、推理路径 & verifier、采样、搜索、reward 设计。 \\
Agent 后训练 & 多步环境任务、工具链、端到端完成率 & rollout infra、任务环境、工具编排、成本与安全评估。 \\
\bottomrule
\end{tabularx}

\subsection{RL scaling 的难点}

Agent 上怎么做好 RL scaling，是罗福莉认为当前清晰但仍需探索的方向。困难在于，Agent 任务的 reward 常常不如代码测试或数学答案清晰；任务环境更长、更开放、更容易出隐藏状态；成本也更高。RL scaling 不只是放大采样，而是要让任务、环境、反馈、成本和安全都能规模化。

\begin{importantbox}{Agent 后训练的关键约束}
Agent 后训练的瓶颈不是单纯缺算法，而是缺可扩展的任务环境和反馈机制。没有稳定环境，rollout 没法规模化；没有可靠 verifier，RL 会学偏；没有成本控制，系统无法生产化。
\end{importantbox}

\subsection{本章小结}

Agent 时代的后训练是系统工程。它要求团队从模型训练视角转向环境闭环视角，把 rollout、reward、policy、evaluation、scaffold 和成本控制连成一套可扩展基础设施。

